% Phần M1 cho Chương 3 — được \input bên trong chapter3_ir.tex nên KHÔNG có \chapter.
% Chỉ dùng các gói đã khai báo trong main.tex; không thêm gói mới.
\section{Thiết kế tầng ngôn ngữ M1}
\label{sec:m1-thiet-ke}

Tầng ngôn ngữ M1 biến văn bản thô thành dãy token sạch để M2 lập chỉ mục. Tầng này
gồm ba khối độc lập: chuẩn hóa chuỗi, tách từ theo so khớp cực đại trên từ điển, và
lọc từ dừng tầng~1 kèm danh sách bảo vệ. Toàn bộ tầng hoạt động ngoại tuyến, không
gọi dịch vụ ngoài, và mọi tài nguyên mà nó dùng đều được ghim bằng hàm băm.

\subsection{Chuẩn hóa văn bản và bảo toàn ranh giới}

Hàm chuẩn hóa \texttt{normalize(text)} trả về dạng hiển thị: chuẩn hóa Unicode
NFC, hạ chữ thường, loại bỏ thẻ đánh dấu, đường dẫn và địa chỉ thư điện tử, giữ
lại chữ cái tiếng Việt cùng khoảng trắng, và thu gọn mọi khoảng trắng về một dấu
cách. Một ví dụ:

\begin{quote}
\texttt{"Tri thức, nhân loại! 2026"} $\rightarrow$ \texttt{"tri thức nhân loại"}.
\end{quote}

Một quyết định thiết kế quan trọng là \emph{dấu câu là ranh giới, không phải
nhiễu}. Dạng hiển thị gộp khoảng trắng nên trông như một chuỗi liên tục, nhưng bộ
tách từ phải không bao giờ ghép một từ xuyên qua dấu câu. Vì lý do đó, bộ tách từ
tiêu thụ một dạng trung gian giữ lại ranh giới, ở đó đoạn văn bản bị cắt tại mỗi
dấu câu và tại mỗi khoảng trắng không phải dấu cách thường (xuống dòng, tab). Các
khoảng được gỡ bỏ (đường dẫn, thư điện tử, thẻ đánh dấu) cũng được thay bằng một
ranh giới cứng chứ không phải bằng dấu cách; nếu thay bằng dấu cách thì hai mảnh
văn bản hai bên sẽ bị dính thành một từ, ví dụ \texttt{"học<br>sinh"} phải tách
thành hai âm tiết chứ không được ghép thành một từ.

Cũng trong khối này có một chỉnh sửa so với mẫu biểu thức chính quy ban đầu. Mẫu
cũ chỉ giữ chữ cái tiếng Việt nên đã vô tình xóa mọi chữ cái Latin, biến
``Windows'' thành ``indows'' và ``Java'' thành ``ava''. Bốn chữ cái Latin
\texttt{f}, \texttt{j}, \texttt{w}, \texttt{z} được giữ lại vì chúng phân biệt
trong các tên vay mượn phổ biến (Windows, WWW, Java, Zalo). Chúng \emph{không}
được thêm vào từ điển; kết quả trung thực là chúng trở thành token đơn ngoài từ
điển thay vì bị bóp méo thành từ sai.

\subsection{Từ điển và cây Trie}

Tách từ dựa trên từ điển Viet74K \cite{viet74k} gồm 69.858 mục, được ghim hàm băm
SHA256 \texttt{175e8b1b...d8ab} và đi kèm thông tin giấy phép cùng nguồn gốc trong
kho. Từ điển được nạp vào một cây Trie hỗ trợ truy vấn theo tiền tố, cho phép bộ
tách từ kiểm tra đoạn ứng viên có trong từ điển hay không mà không phải so khớp
toàn bộ danh sách. Cửa sổ trượt mặc định lấy bằng độ dài mục dài nhất trong từ
điển, nên không bỏ sót từ nào; người dùng vẫn có thể chỉ định một giới hạn cửa sổ
nhỏ hơn khi cần quan sát thuật toán.

\subsection{Thuật toán FMM, BMM và BiMM}

Ba phương pháp dùng chung một cửa sổ trượt nhưng khác hướng duyệt. FMM duyệt tiến,
BMM duyệt lùi. BiMM chạy cả hai rồi chọn theo bộ quy tắc phân xử sau, xét theo
thứ tự:

\begin{enumerate}
  \item ít token hơn;
  \item ít token một âm tiết hơn;
  \item ít token ngoài từ điển hơn;
  \item nếu vẫn hòa, chọn BMM.
\end{enumerate}

Tiêu chí thứ ba là lưới an toàn tường minh: với cài đặt Trie thuần, mọi token
ngoài từ điển đều được phát ra dưới dạng một âm tiết, nên tiêu chí này được suy ra
từ tiêu chí thứ hai; nhóm vẫn giữ nó cùng khóa so sánh đầy đủ để người đọc kiểm
chứng lại quyết định thay vì tin vào phần tóm tắt.

Ví dụ nhập nhằng kinh điển minh họa vì sao cần phân xử. Với câu ``học sinh học
sinh học'' và một từ điển fixture có cả ``học sinh'' lẫn ``sinh học'':

\begin{quote}
FMM: \texttt{học\_sinh, học\_sinh, học}\\
BMM: \texttt{học, sinh\_học, sinh\_học}\\
BiMM: \texttt{học, sinh\_học, sinh\_học}
\end{quote}

Hai phương án có cùng số token và cùng số token một âm tiết, nên hòa và BiMM chọn
BMM. Ví dụ này dùng một từ điển fixture minh bạch, không phụ thuộc từ điển đóng
gói, để phép so sánh là một khẳng định về \emph{thuật toán} chứ không về một từ
điển cụ thể.

Mỗi lần chạy đều sinh một \emph{vết} từng bước. Mỗi dòng vết ghi số thứ tự bước,
hướng duyệt, chỉ số âm tiết bắt đầu và kết thúc (kết thúc mở), đoạn ứng viên, độ
dài cửa sổ, cờ có khớp từ điển hay không, và hành động -- một trong
\texttt{MATCHED}, \texttt{SHRINK} hoặc \texttt{SINGLE}. Vết dùng đúng từ điển và
cửa sổ của lần chạy, nên người đọc có thể tái lập từng quyết định thay vì chỉ thấy
kết quả cuối cùng. Với đầu vào rỗng hoặc chỉ gồm dấu câu, kết quả là dãy token rỗng
và vết rỗng.

\subsection{Lọc từ dừng tầng 1 và danh sách bảo vệ}

Khối thứ ba loại các hư từ ngữ pháp -- đại từ, giới từ, trợ từ, phó từ chỉ mức độ,
liên từ -- vốn đóng góp ít cho phân biệt truy xuất. Hai tài nguyên có nguồn và quy
tắc biên soạn điều khiển bước này: danh sách hư từ ngữ pháp và danh sách bảo vệ
(guard). Tập từ dừng hiệu lực là hiệu của hai danh sách: \emph{guard luôn thắng}.

Quy tắc bảo vệ có phạm vi rộng hơn một so khớp đơn giản: một token được bảo vệ nếu
chính nó là guard, hoặc nếu một thành phần của nó (tách theo dấu gạch nối âm tiết)
cũng là guard. Nhờ đó việc bảo vệ ``không'' cũng giữ nguyên ``không\_thể'', bảo vệ
``vì'' giữ nguyên ``vì\_dầu'', bảo vệ ``chẳng'' giữ nguyên ``chẳng\_hạn''. Điểm
tinh tế là chỉ những thành phần \emph{tự chúng là guard} mới có thể bảo vệ: nếu mọi
âm tiết của mọi guard đều có thể bảo vệ thì ``thể'' sẽ được bảo vệ và ``thể\_thao''
bị bảo vệ nhầm. Chính vì vậy các từ ghép cần bảo vệ vì lý do khác được liệt kê
tường minh trong tệp guard.

Ví dụ minh họa:

\begin{quote}
\texttt{filter\_tier1(["bóng", "đá", "và", "của", "bóng\_đá"])}
$\rightarrow$ \texttt{["bóng", "đá", "bóng\_đá"]}
\end{quote}

Ở đây ``và'' và ``của'' bị loại, còn ``bóng\_đá'' là một từ ghép nội dung được giữ
nguyên. Cùng cơ chế bảo vệ này được tái dùng cho tầng~2 khi TDV chọn stoplist, để
các từ phủ định không bị xóa ở cả hai tầng.

\subsection{Tài nguyên và xuất xứ}

Ba tài nguyên ngôn ngữ học được đóng gói kèm số mục thực tế và hàm băm, để kết quả
tái lập được và để bất kỳ ai cũng kiểm tra lại được dữ liệu đầu vào.

\begin{table}[ht]
\centering
\begin{tabular}{|p{0.30\textwidth}|p{0.38\textwidth}|p{0.18\textwidth}|}
\hline
Tài nguyên & Vai trò & Số mục \\
\hline
\texttt{dict.txt} & Từ điển tách từ (Viet74K) & 69.858 \\
\texttt{grammar\_\allowbreak stopwords.txt} & Hư từ ngữ pháp tầng 1 & 181 \\
\texttt{guard\_\allowbreak whitelist.txt} & Danh sách bảo vệ (guard) & 110 \\
\hline
\end{tabular}
\caption{Ba tài nguyên của tầng ngôn ngữ M1 (đặt trong thư mục \texttt{resources/})
và số mục tại thời điểm viết. Số mục được đọc trực tiếp từ tệp; danh sách hư từ sau
khi trừ guard còn 124 từ dừng hiệu lực.}
\label{tab:m1-tai-nguyen}
\end{table}

Từ điển là danh sách mã nguồn mở có ghi giấy phép \cite{viet74k}. Danh sách từ
dừng tĩnh dùng cho cấu hình C2 là một bản chụp có nguồn, phiên bản và thông báo
giấy phép MIT \cite{static_stopwords}. Hai danh sách tầng~1 do nhóm biên soạn kèm
quy tắc và số mục thực tế; các từ dừng được gỡ lại chỉ là các từ phủ định và logic
có lý do ngữ nghĩa rõ ràng, không bù danh sách bằng thực từ để đạt một con số.

\subsection{Kiểm chứng tầng ngôn ngữ}

Ba phương pháp tách từ được kiểm chứng bằng từ điển fixture minh bạch và bằng một
đường tham chiếu trực tiếp trên từ điển/câu ngẫu nhiên. Vết thuật toán cho biết ứng
viên và quyết định tại từng bước; dãy token do \texttt{tokenize} trả về được kiểm
riêng bằng fixture, vì với BMM các dòng ứng viên được đánh dấu có thể chồng lấn nên
không dùng chúng để tái dựng kết quả. Bộ lọc tầng~1 được kiểm chứng
bằng các ca guard: phủ định và từ ghép chứa guard phải được giữ nguyên, còn hư từ
không nằm trong guard phải bị loại. Các ca kiểm thử này chứng minh \emph{luật hoạt
động đúng}, không đo \emph{chất lượng} tách từ.

Chất lượng tách từ chỉ đo được khi có nhãn người. Nhóm đã dựng bộ mẫu 150 câu
(15 câu mỗi lĩnh vực) rút từ kho ngữ liệu đã khóa, kèm công cụ gán nhãn ngoại
tuyến và bộ đo khớp span từ. Hiện 150/150 câu có cờ xác nhận và phép đo word-span
đã chạy, nhưng ba câu SG010, SG053, SG089 còn cần rà lại về ranh giới từ; do đó
F1 ở Chương~4 là số đo tạm thời trên gold hiện có, không phải kết luận chất lượng
cuối. Gợi ý của máy chỉ phục vụ gán nhãn nhanh và không tự động trở thành gold.
Quy trình đo, các ràng buộc theo hàm băm và giới hạn của phép đo trình bày trong
Chương~4.
